% Options for packages loaded elsewhere
\PassOptionsToPackage{unicode}{hyperref}
\PassOptionsToPackage{hyphens}{url}
\documentclass[
]{article}
\usepackage{xcolor}
\usepackage[margin=1in]{geometry}
\usepackage{amsmath,amssymb}
\setcounter{secnumdepth}{-\maxdimen} % remove section numbering
\usepackage{iftex}
\ifPDFTeX
  \usepackage[T1]{fontenc}
  \usepackage[utf8]{inputenc}
  \usepackage{textcomp} % provide euro and other symbols
\else % if luatex or xetex
  \usepackage{unicode-math} % this also loads fontspec
  \defaultfontfeatures{Scale=MatchLowercase}
  \defaultfontfeatures[\rmfamily]{Ligatures=TeX,Scale=1}
\fi
\usepackage{lmodern}
\ifPDFTeX\else
  % xetex/luatex font selection
\fi
% Use upquote if available, for straight quotes in verbatim environments
\IfFileExists{upquote.sty}{\usepackage{upquote}}{}
\IfFileExists{microtype.sty}{% use microtype if available
  \usepackage[]{microtype}
  \UseMicrotypeSet[protrusion]{basicmath} % disable protrusion for tt fonts
}{}
\makeatletter
\@ifundefined{KOMAClassName}{% if non-KOMA class
  \IfFileExists{parskip.sty}{%
    \usepackage{parskip}
  }{% else
    \setlength{\parindent}{0pt}
    \setlength{\parskip}{6pt plus 2pt minus 1pt}}
}{% if KOMA class
  \KOMAoptions{parskip=half}}
\makeatother
\usepackage{color}
\usepackage{fancyvrb}
\newcommand{\VerbBar}{|}
\newcommand{\VERB}{\Verb[commandchars=\\\{\}]}
\DefineVerbatimEnvironment{Highlighting}{Verbatim}{commandchars=\\\{\}}
% Add ',fontsize=\small' for more characters per line
\usepackage{framed}
\definecolor{shadecolor}{RGB}{248,248,248}
\newenvironment{Shaded}{\begin{snugshade}}{\end{snugshade}}
\newcommand{\AlertTok}[1]{\textcolor[rgb]{0.94,0.16,0.16}{#1}}
\newcommand{\AnnotationTok}[1]{\textcolor[rgb]{0.56,0.35,0.01}{\textbf{\textit{#1}}}}
\newcommand{\AttributeTok}[1]{\textcolor[rgb]{0.13,0.29,0.53}{#1}}
\newcommand{\BaseNTok}[1]{\textcolor[rgb]{0.00,0.00,0.81}{#1}}
\newcommand{\BuiltInTok}[1]{#1}
\newcommand{\CharTok}[1]{\textcolor[rgb]{0.31,0.60,0.02}{#1}}
\newcommand{\CommentTok}[1]{\textcolor[rgb]{0.56,0.35,0.01}{\textit{#1}}}
\newcommand{\CommentVarTok}[1]{\textcolor[rgb]{0.56,0.35,0.01}{\textbf{\textit{#1}}}}
\newcommand{\ConstantTok}[1]{\textcolor[rgb]{0.56,0.35,0.01}{#1}}
\newcommand{\ControlFlowTok}[1]{\textcolor[rgb]{0.13,0.29,0.53}{\textbf{#1}}}
\newcommand{\DataTypeTok}[1]{\textcolor[rgb]{0.13,0.29,0.53}{#1}}
\newcommand{\DecValTok}[1]{\textcolor[rgb]{0.00,0.00,0.81}{#1}}
\newcommand{\DocumentationTok}[1]{\textcolor[rgb]{0.56,0.35,0.01}{\textbf{\textit{#1}}}}
\newcommand{\ErrorTok}[1]{\textcolor[rgb]{0.64,0.00,0.00}{\textbf{#1}}}
\newcommand{\ExtensionTok}[1]{#1}
\newcommand{\FloatTok}[1]{\textcolor[rgb]{0.00,0.00,0.81}{#1}}
\newcommand{\FunctionTok}[1]{\textcolor[rgb]{0.13,0.29,0.53}{\textbf{#1}}}
\newcommand{\ImportTok}[1]{#1}
\newcommand{\InformationTok}[1]{\textcolor[rgb]{0.56,0.35,0.01}{\textbf{\textit{#1}}}}
\newcommand{\KeywordTok}[1]{\textcolor[rgb]{0.13,0.29,0.53}{\textbf{#1}}}
\newcommand{\NormalTok}[1]{#1}
\newcommand{\OperatorTok}[1]{\textcolor[rgb]{0.81,0.36,0.00}{\textbf{#1}}}
\newcommand{\OtherTok}[1]{\textcolor[rgb]{0.56,0.35,0.01}{#1}}
\newcommand{\PreprocessorTok}[1]{\textcolor[rgb]{0.56,0.35,0.01}{\textit{#1}}}
\newcommand{\RegionMarkerTok}[1]{#1}
\newcommand{\SpecialCharTok}[1]{\textcolor[rgb]{0.81,0.36,0.00}{\textbf{#1}}}
\newcommand{\SpecialStringTok}[1]{\textcolor[rgb]{0.31,0.60,0.02}{#1}}
\newcommand{\StringTok}[1]{\textcolor[rgb]{0.31,0.60,0.02}{#1}}
\newcommand{\VariableTok}[1]{\textcolor[rgb]{0.00,0.00,0.00}{#1}}
\newcommand{\VerbatimStringTok}[1]{\textcolor[rgb]{0.31,0.60,0.02}{#1}}
\newcommand{\WarningTok}[1]{\textcolor[rgb]{0.56,0.35,0.01}{\textbf{\textit{#1}}}}
\usepackage{graphicx}
\makeatletter
\newsavebox\pandoc@box
\newcommand*\pandocbounded[1]{% scales image to fit in text height/width
  \sbox\pandoc@box{#1}%
  \Gscale@div\@tempa{\textheight}{\dimexpr\ht\pandoc@box+\dp\pandoc@box\relax}%
  \Gscale@div\@tempb{\linewidth}{\wd\pandoc@box}%
  \ifdim\@tempb\p@<\@tempa\p@\let\@tempa\@tempb\fi% select the smaller of both
  \ifdim\@tempa\p@<\p@\scalebox{\@tempa}{\usebox\pandoc@box}%
  \else\usebox{\pandoc@box}%
  \fi%
}
% Set default figure placement to htbp
\def\fps@figure{htbp}
\makeatother
\setlength{\emergencystretch}{3em} % prevent overfull lines
\providecommand{\tightlist}{%
  \setlength{\itemsep}{0pt}\setlength{\parskip}{0pt}}
\usepackage{bookmark}
\IfFileExists{xurl.sty}{\usepackage{xurl}}{} % add URL line breaks if available
\urlstyle{same}
\hypersetup{
  pdftitle={Bayesian ANOVA},
  hidelinks,
  pdfcreator={LaTeX via pandoc}}

\title{Bayesian ANOVA}
\author{}
\date{\vspace{-2.5em}2026-04-17}

\begin{document}
\maketitle

\subsection{Introduction}\label{introduction}

Classical ANOVA reports an \(F\)-statistic, a \(p\)-value, and --- if
the omnibus is significant --- pairwise comparisons with a
multiple-testing correction tacked on. A Bayesian ANOVA replaces all of
that with a single posterior distribution over the group means and any
contrast you care to compute. There is no omnibus to clear before
contrasts become legitimate, and there is no multiple-comparison
adjustment to argue about, because every contrast already has its own
honest credible interval reflecting the joint posterior of the
parameters that compose it.

\subsection{Prerequisites}\label{prerequisites}

A working understanding of one-way ANOVA, the distinction between fixed
and random effects, and Bayesian regression. Familiarity with
\texttt{emmeans} for marginal means makes the contrast step natural.

\subsection{Theory}\label{theory}

The model is \(y_{ij} \sim \mathcal N(\mu_i, \sigma^2)\) with
\(\mu_i = \mu + \alpha_i\). Two parameterisations are common: a
fixed-effects parameterisation with weakly informative priors on the
contrasts, and a hierarchical parameterisation with
\(\alpha_i \sim \mathcal N(0, \tau^2)\) that pools group estimates
toward the grand mean. The hierarchical version is preferable when group
counts are large or when groups represent samples from a broader
population. Pairwise contrasts and any linear combination of \(\mu_i\)
are computed at the level of MCMC draws, so each draw produces a
contrast value and the posterior of the contrast is the empirical
distribution of those values.

\subsection{Assumptions}\label{assumptions}

Approximately Normal residuals within groups, homogeneity of variance
unless \texttt{sigma} is also modelled by group, exchangeability of
groups under the hierarchical version, and proper priors. Independence
between observations after group is the same Bayesian assumption as in
classical ANOVA.

\subsection{R Implementation}\label{r-implementation}

\begin{Shaded}
\begin{Highlighting}[]
\FunctionTok{library}\NormalTok{(brms); }\FunctionTok{library}\NormalTok{(bayestestR)}

\FunctionTok{data}\NormalTok{(PlantGrowth)}

\NormalTok{fit }\OtherTok{\textless{}{-}} \FunctionTok{brm}\NormalTok{(weight }\SpecialCharTok{\textasciitilde{}}\NormalTok{ group, }\AttributeTok{data =}\NormalTok{ PlantGrowth,}
           \AttributeTok{prior =} \FunctionTok{prior}\NormalTok{(}\FunctionTok{normal}\NormalTok{(}\DecValTok{0}\NormalTok{, }\DecValTok{1}\NormalTok{), }\AttributeTok{class =} \StringTok{"b"}\NormalTok{),}
           \AttributeTok{chains =} \DecValTok{4}\NormalTok{, }\AttributeTok{iter =} \DecValTok{2000}\NormalTok{, }\AttributeTok{refresh =} \DecValTok{0}\NormalTok{)}

\CommentTok{\# Pairwise contrasts}
\NormalTok{pairs\_summary }\OtherTok{\textless{}{-}} \FunctionTok{pairs}\NormalTok{(emmeans}\SpecialCharTok{::}\FunctionTok{emmeans}\NormalTok{(fit, }\SpecialCharTok{\textasciitilde{}}\NormalTok{ group))}
\NormalTok{pairs\_summary}

\CommentTok{\# Bayes factors for effect}
\FunctionTok{bayesfactor\_parameters}\NormalTok{(fit)}
\end{Highlighting}
\end{Shaded}

\subsection{Output \& Results}\label{output-results}

\texttt{emmeans()} returns the posterior marginal mean for each group;
\texttt{pairs()} returns the posterior of every pairwise difference with
an HPD or equal-tailed credible interval.
\texttt{bayesfactor\_parameters()} evaluates the Savage-Dickey density
ratio at zero for each contrast, returning a Bayes factor for ``non-zero
contrast'' against the null.

\subsection{Interpretation}\label{interpretation}

A reporting sentence: ``Posterior contrasts from a Bayesian one-way
ANOVA showed \texttt{trt1} \(-\) \texttt{ctrl} \(= -0.37\) (95 \%
credible interval \(-0.95\) to \(0.21\), \(\mathrm{BF}_{10} = 1.1\)),
\texttt{trt2} \(-\) \texttt{ctrl} \(= 0.49\) (95 \% CrI 0.04 to 0.94,
\(\mathrm{BF}_{10} = 4.6\)); the data weakly support a treatment-2
effect and are inconclusive about treatment 1.'' Pair the credible
interval with the Bayes factor only when the prior under the alternative
is disclosed, since the latter depends sensitively on it.

\subsection{Practical Tips}\label{practical-tips}

\begin{itemize}
\tightlist
\item
  For more than five or six groups, switch to a hierarchical
  \((1 \mid \mathrm{group})\) random-effects parameterisation; shrinkage
  stabilises small-group estimates and gives a principled
  population-level summary.
\item
  \texttt{emmeans()} works directly on \texttt{brms} fits and integrates
  with \texttt{pairs()}, \texttt{contrast()}, and \texttt{confint()} to
  produce posterior contrasts on whatever scale you need.
\item
  Credible intervals already adjust for the joint posterior;
  multiple-comparison corrections in the frequentist sense are
  unnecessary, but you can still apply a region-of-practical-equivalence
  (ROPE) test if you want a decision rule.
\item
  For unequal variances, fit
  \texttt{bf(weight\ \textasciitilde{}\ group,\ sigma\ \textasciitilde{}\ group)}
  so the residual scale is also estimated per group; the contrast
  machinery generalises without changes.
\item
  Report posteriors graphically (\texttt{mcmc\_areas},
  \texttt{stat\_halfeye}) alongside the numeric summaries; readers
  absorb the joint structure faster from a plot than from a contrast
  table.
\end{itemize}

\subsection{R Packages Used}\label{r-packages-used}

\texttt{brms} for model specification and sampling, \texttt{emmeans} for
posterior marginal means and contrasts, \texttt{bayestestR} for Bayes
factors and ROPE summaries, and \texttt{tidybayes} for tidy posterior
draws of contrasts when custom plots are needed.

\subsection{For Reviewers}\label{for-reviewers}

\textbf{What to look for in a paper using this method.}

\begin{itemize}
\tightlist
\item
  \textbf{Common misapplications.}
\item
  \textbf{Diagnostics that should be reported but often aren't.}
\item
  \textbf{Red flags in tables and figures.}
\item
  \textbf{What to verify.}
\item
  \textbf{What an adequate Methods paragraph must contain.}
\end{itemize}

\subsection{See also --- labs in this
chapter}\label{see-also-labs-in-this-chapter}

\begin{itemize}
\tightlist
\item
  \href{labs/lab-bayesian-thinking-control-vs-decision-error.Rmd}{Bayesian
  thinking; α control vs decision error}
\item
  \href{labs/lab-brms-stan-loo-hierarchical-models.Rmd}{brms / Stan,
  LOO, hierarchical models}
\end{itemize}

\end{document}
